\documentclass[12pt,a4paper]{report}

% ===================== PACKAGES =====================
\usepackage[utf8]{inputenc}
\usepackage[T1]{fontenc}
\usepackage[french]{babel}
\usepackage{graphicx}
\usepackage{geometry}
\usepackage{xcolor}
\usepackage{hyperref}
\usepackage{listings}
\usepackage{tikz}
\usepackage{pgf-umlsd}   % diagrammes de sequence UML
\usepackage{array}
\usepackage{booktabs}
\usepackage{titlesec}

\geometry{margin=2.5cm}
\hypersetup{colorlinks=true, linkcolor=blue, urlcolor=blue}

% Couleurs pour le code
\definecolor{codegreen}{rgb}{0,0.6,0}
\definecolor{codegray}{rgb}{0.5,0.5,0.5}
\definecolor{backcolour}{rgb}{0.95,0.95,0.92}
\lstset{
  backgroundcolor=\color{backcolour},
  commentstyle=\color{codegreen},
  keywordstyle=\color{blue},
  numberstyle=\tiny\color{codegray},
  basicstyle=\ttfamily\footnotesize,
  breaklines=true, numbers=left, frame=single, captionpos=b
}

% ===================== DOCUMENT =====================
\begin{document}

% ---------- PAGE DE GARDE ----------
\begin{titlepage}
\centering
\vspace*{1cm}
{\Large\textbf{Université SESAME}}\\[0.3cm]
{\large Année universitaire 2025--2026}\\[2cm]

\rule{\linewidth}{0.5mm}\\[0.5cm]
{\Huge\bfseries LoL Fantasy League}\\[0.4cm]
{\Large Plateforme web de fantasy esports\\League of Legends}\\[0.3cm]
\rule{\linewidth}{0.5mm}\\[2cm]

{\large \textbf{Projet de fin d'année}}\\[2cm]

\begin{flushleft}
\large
\textbf{Réalisé par :} Oumaima Kammoun\\[0.3cm]
\textbf{Encadré par :} [Nom de l'encadrant]\\[0.3cm]
\textbf{Filière :} [Votre filière]\\
\end{flushleft}

\vfill
{\large \today}
\end{titlepage}

% ---------- REMERCIEMENTS ----------
\chapter*{Remerciements}
\addcontentsline{toc}{chapter}{Remerciements}
Je tiens à remercier mon encadrant pour son suivi et ses conseils tout au long
de ce projet, ainsi que le corps enseignant de l'Université SESAME pour la
qualité de la formation dispensée.

% ---------- RESUME ----------
\chapter*{Résumé}
\addcontentsline{toc}{chapter}{Résumé}
Ce projet consiste en la conception et la réalisation d'une plateforme web de
\textit{fantasy esports} dédiée à League of Legends. L'utilisateur compose une
équipe virtuelle de cinq joueurs professionnels réels et gagne des points selon
leurs performances lors des matchs officiels. La plateforme intègre l'API
officielle LoL Esports en temps réel et un moteur de scoring entièrement
automatique. Elle repose sur une architecture trois tiers : un frontend React,
une API REST Django et une base PostgreSQL.\\[0.3cm]
\textbf{Mots-clés :} Django, React, API REST, fantasy esports, JWT, 2FA, temps réel.

% ---------- TABLE DES MATIERES ----------
\tableofcontents

% =====================================================
\chapter{Introduction générale}
% =====================================================
Le \textit{fantasy sport} est un type de jeu où les participants constituent des
équipes virtuelles de joueurs réels et marquent des points selon leurs
performances effectives. Appliqué à l'esport League of Legends, ce concept
permet de prolonger l'expérience des compétitions professionnelles.

\section{Contexte}
Les compétitions League of Legends (LCK, LEC, LCS, LPL, EMEA Masters) rassemblent
des millions de spectateurs. Ce projet propose une plateforme permettant aux
fans de participer activement en composant leurs équipes fantasy.

\section{Problématique}
Comment concevoir une plateforme qui : (1) exploite des données réelles et à jour
des matchs professionnels, (2) calcule automatiquement et de façon fiable les
points, et (3) reste fonctionnelle en permanence sans intervention manuelle ?

\section{Objectifs}
\begin{itemize}
  \item Concevoir une application full-stack sécurisée (JWT + 2FA).
  \item Intégrer l'API officielle LoL Esports en temps réel.
  \item Développer un moteur de scoring automatique et fiable.
  \item Offrir une expérience complète : ligues, rosters, marché, social, chatbot.
\end{itemize}

% =====================================================
\chapter{Analyse et spécification des besoins}
% =====================================================
\section{Besoins fonctionnels}
\begin{itemize}
  \item \textbf{Authentification} : inscription, connexion à deux facteurs (2FA).
  \item \textbf{Ligues} : création/adhésion via code d'invitation.
  \item \textbf{Roster} : composition d'une équipe de 5 joueurs avec budget et capitaine.
  \item \textbf{Scoring} : calcul automatique des points à chaque match.
  \item \textbf{Marché} : historique des achats et ventes de joueurs.
  \item \textbf{Social} : classements, pronostics, abonnements.
  \item \textbf{Chatbot} : assistant intelligent.
\end{itemize}

\section{Besoins non fonctionnels}
\begin{itemize}
  \item \textbf{Sécurité} : JWT, 2FA, contrôle d'accès par rôle (RBAC).
  \item \textbf{Fiabilité} : scoring idempotent, fonctionnement automatique permanent.
  \item \textbf{Temps réel} : matchs en direct, scores actualisés.
  \item \textbf{Données réelles} : aucune donnée fictive, API officielle.
\end{itemize}

\section{Acteurs du système}
\begin{itemize}
  \item \textbf{Joueur} : participe, compose son roster, fait des pronostics.
  \item \textbf{Manager} : crée et gère ses ligues privées (hérite du joueur).
  \item \textbf{Administrateur} : gère les joueurs pros et les comptes.
\end{itemize}

% =====================================================
\chapter{Conception}
% =====================================================
\section{Architecture générale}
L'application suit une architecture \textbf{trois tiers} :
\begin{itemize}
  \item \textbf{Couche présentation} : application React 19 (interface utilisateur).
  \item \textbf{Couche logique} : API REST Django + Django REST Framework.
  \item \textbf{Couche données} : base PostgreSQL.
\end{itemize}
Une API externe (LoL Esports) alimente les données des matchs en temps réel.

\section{Diagrammes de séquence}

% -------- SEQUENCE 1 : Authentification 2FA --------
\subsection{Authentification avec double facteur (2FA)}
\begin{figure}[h]
\centering
\resizebox{0.95\textwidth}{!}{
\begin{sequencediagram}
\newthread{u}{Utilisateur}
\newinst{f}{Frontend React}
\newinst{a}{API Django}
\newinst{d}{Base PostgreSQL}

\begin{call}{u}{Saisie email + mot de passe}{f}{}
\end{call}
\begin{call}{f}{POST /auth/login/}{a}{user\_id, otp\_method}
  \begin{call}{a}{Vérifie identifiants}{d}{OK}
  \end{call}
\end{call}
\begin{call}{u}{Choix méthode OTP}{f}{}
\end{call}
\begin{call}{f}{POST /auth/login/choose-otp/}{a}{QR code / code email}
\end{call}
\begin{call}{u}{Saisie code OTP}{f}{}
\end{call}
\begin{call}{f}{POST /auth/login/verify/}{a}{tokens JWT}
  \begin{call}{a}{Vérifie code OTP}{d}{valide}
  \end{call}
\end{call}
\begin{call}{f}{Stocke access + refresh token}{f}{}
\end{call}
\end{sequencediagram}
}
\caption{Séquence d'authentification à deux facteurs}
\end{figure}

% -------- SEQUENCE 2 : Composition roster --------
\subsection{Composition d'un roster et achat de joueur}
\begin{figure}[h]
\centering
\resizebox{0.95\textwidth}{!}{
\begin{sequencediagram}
\newthread{u}{Joueur}
\newinst{f}{Frontend}
\newinst{a}{API Django}
\newinst{d}{Base PostgreSQL}

\begin{call}{u}{Sélectionne un joueur}{f}{}
\end{call}
\begin{call}{f}{POST /rosters/<id>/add/}{a}{roster mis à jour}
  \begin{call}{a}{Vérifie budget suffisant}{d}{OK}
  \end{call}
  \begin{call}{a}{Crée RosterSlot}{d}{}
  \end{call}
  \begin{call}{a}{Crée Transfer (achat)}{d}{}
  \end{call}
\end{call}
\begin{call}{u}{Désigne le capitaine}{f}{}
\end{call}
\begin{call}{f}{POST /rosters/<id>/captain/}{a}{OK}
  \begin{call}{a}{Met is\_captain=True}{d}{}
  \end{call}
\end{call}
\end{sequencediagram}
}
\caption{Séquence de composition d'un roster}
\end{figure}

% -------- SEQUENCE 3 : Scoring automatique --------
\subsection{Calcul automatique des scores}
\begin{figure}[h]
\centering
\resizebox{0.95\textwidth}{!}{
\begin{sequencediagram}
\newthread{s}{Scheduler (3 min)}
\newinst{a}{API Django}
\newinst{e}{API LoL Esports}
\newinst{d}{Base PostgreSQL}

\begin{call}{s}{Déclenche la synchronisation}{a}{}
  \begin{call}{a}{get\_schedule()}{e}{matchs terminés}
  \end{call}
  \begin{call}{a}{Récupère stats joueurs}{e}{stats (ou vide)}
  \end{call}
  \begin{call}{a}{Si pas de stats : génère stats de secours}{d}{}
  \end{call}
  \begin{call}{a}{calculate\_scores\_for\_match()}{d}{FantasyScore créés}
  \end{call}
  \begin{call}{a}{Recalcule totaux (idempotent)}{d}{classements à jour}
  \end{call}
  \begin{call}{a}{Valide les pronostics}{d}{+5 pts si correct}
  \end{call}
\end{call}
\end{sequencediagram}
}
\caption{Séquence du calcul automatique des scores}
\end{figure}

\section{Modèle de données}
Les principales entités sont : \texttt{Utilisateur}, \texttt{League},
\texttt{LeagueMember}, \texttt{Roster}, \texttt{RosterSlot}, \texttt{Player},
\texttt{Match}, \texttt{MatchPlayerStat}, \texttt{FantasyScore},
\texttt{Transfer}, \texttt{Pronostic}. Un \texttt{Roster} appartient à un
\texttt{Utilisateur} et une \texttt{League}, et contient 5 \texttt{RosterSlot}
liés chacun à un \texttt{Player}.

% =====================================================
\chapter{Réalisation}
% =====================================================
\section{Technologies utilisées}
\begin{center}
\begin{tabular}{ll}
\toprule
\textbf{Composant} & \textbf{Technologie} \\
\midrule
Frontend & React 19, react-router-dom, Axios \\
Backend & Django 6, Django REST Framework \\
Authentification & SimpleJWT + 2FA (TOTP / OTP email) \\
Base de données & PostgreSQL 18 \\
API externe & LoL Esports API \\
Chatbot & Ollama (llama3.2) \\
\bottomrule
\end{tabular}
\end{center}

\section{Moteur de scoring (KPI)}
La formule de calcul des points combine performance individuelle et collective :
\[
\text{KPI} = \left( \sum S_{\text{indiv}} \times 0.70 \right) + \left( S_{\text{équipe}} \times 0.30 \right)
\]
avec
\[
S_{\text{indiv}} = K\times 3 + A\times 2 + D\times(-1.5) + CS\times 0.02 + VS\times 0.1
\]
\[
S_{\text{collectif}} = \text{Victoire}\times 5 + \text{Barons}\times 2 + \text{Dragons}\times 2 + \text{Tours}\times 1
\]
Le capitaine voit son score individuel multiplié par 1,5.

\section{Système automatique}
Un \textit{scheduler} en arrière-plan synchronise les matchs toutes les 3 minutes.
Le calcul des scores est \textbf{idempotent} : les totaux sont recalculés depuis
la somme des scores, ce qui garantit l'absence de double comptage même en cas de
ré-exécution. Lorsque l'API ne fournit pas les statistiques détaillées, le système
génère des statistiques de secours cohérentes afin que les points soient
toujours calculés.

% =====================================================
\chapter{Difficultés rencontrées et solutions}
% =====================================================
\begin{itemize}
  \item \textbf{Statistiques manquantes} : l'API ne fournit pas toujours les stats
        détaillées. \textit{Solution :} génération de stats de secours.
  \item \textbf{Double comptage} : le scoring incrémentait les totaux.
        \textit{Solution :} rendre le calcul idempotent.
  \item \textbf{Matchs futurs marqués terminés} : données obsolètes de l'API.
        \textit{Solution :} garde temporelle (un match futur n'est jamais terminé).
  \item \textbf{Points rétroactifs} : matchs antérieurs au roster comptés.
        \textit{Solution :} ne compter que les matchs postérieurs à la création du roster.
  \item \textbf{Streams manquants en direct} : \textit{Solution :} streams officiels
        par défaut par ligue.
\end{itemize}

% =====================================================
\chapter{Conclusion et perspectives}
% =====================================================
Ce projet a permis de concevoir une plateforme complète de fantasy esports,
exploitant des données réelles et offrant un scoring entièrement automatique.
L'architecture trois tiers, la sécurité (JWT + 2FA + RBAC) et la résilience du
moteur de scoring constituent les principaux atouts techniques.

\textbf{Perspectives :} système de notifications, badges et récompenses,
statistiques avancées avec graphiques, mode tournoi à élimination directe,
échanges de joueurs entre membres, gestion de saisons.

\end{document}
